\begin{scanpage}{0163}{145}
Now let us remove the auxiliary dipoles to such a distance that the spherical wave arriving from them may be considered to be a plane wave on the section from the half-plane's edge to the point Q. In this case, in accordance with Equation (20.08), the field created by the wave at the point Q will equal
\[
\left.\begin{aligned}E_z(Q)&=E_{0z}(0)[u(d,\varphi'-\varphi'')-u(d,\varphi'+\varphi'')],\\
H_z(Q)&=H_{0z}(0)[u(d,\varphi'-\varphi'')+u(d,\varphi'+\varphi'')].
\end{aligned}\right\}\tag{24.10}
\]
The expressions
\[
E_{0z}(0)=k^2p_{0z}\frac{e^{ikR}}R,\quad H_{0z}(0)=k^2m_{0z}\frac{e^{ikR}}R\tag{24.11}
\]
determine the fields created by the auxiliary dipoles in free space (with the absence of the half-plane) at the point O.

Consequently, the fields excited at the point P by the electric and magnetic dipoles which are found at the point Q above the half-plane equal respectively
\[
\left.\begin{aligned}E_z&=k^2p_z[u(d,\varphi'-\varphi'')-u(d,\varphi'+\varphi'')]\frac{e^{ikR}}R,\\
H_z&=k^2m_z[u(d,\varphi'-\varphi'')+u(d,\varphi'+\varphi'')]\frac{e^{ikR}}R.
\end{aligned}\right\}\tag{24.12}
\]
With the absence of the half-plane, these dipoles create at the point P the field
\[
\left.\begin{aligned}E_z&=k^2p_z\frac{e^{ikR}}R e^{-ikd\cos(\varphi'-\varphi'')},\\
H_z&=k^2m_z\frac{e^{ikR}}R e^{-ikd\cos(\varphi'-\varphi'')}.
\end{aligned}\right\}\tag{24.13}
\]
Comparing Expressions (24.12) and (24.13) we find the shading functions
\[
\left.\begin{aligned}\widetilde F&=[u(d,\varphi'-\varphi'')-u(d,\varphi'+\varphi'')]e^{ikd\cos(\varphi'-\varphi'')},\\
\widetilde G&=[u(d,\varphi'-\varphi'')+u(d,\varphi'+\varphi'')]e^{ikd\cos(\varphi'-\varphi'')}.
\end{aligned}\right\}\tag{24.14}
\]
\end{scanpage}
